\documentclass{article}

%\usepackage[margin=1in]{geometry}

\title{Ch. 27 CW \#1}
\author{Principles of Macroeconomics}
\date{\today}

\begin{document}

\maketitle

%\section*{Instructions}
%All of your homework and classwork for Ch. 27 will come from these problems. Please do them on another sheet of paper and turn them in on Gradescope by the due dates listed on Tigernet.

\section*{27.2. S and I in the National Income Accounts}

\begin{enumerate}

\item If the government collect more in tax revenue than it spends, and households consume more than they get in after-tax income, then

	\begin{enumerate}
	
	\item private saving and public saving are both positive.
	
	\item private saving and public saving are both negative.
	
	\item private saving is positive, but public saving is negative.
	
	\item private saving is negative, but public saving is positive.
	
	\end{enumerate}
	
\item A closed economy has income of \$1,000, government spending of \$200, taxes of \$150, and investment of \$250. What is private saving?

	\begin{enumerate}
	
	\item \$100
	
	\item \$200
	
	\item \$300
	
	\item \$400
	
	\end{enumerate}
	
\item Which of the following situations represent investment and which represent saving? Explain.

	\begin{enumerate}
	
	\item Your family takes out a mortgage and buys a new house.
	
	\item You use your \$200 paycheck to buy stock in AT\&T.
	
	\item Your roommate earns \$100 and deposits it in an account at a bank.
	
	\item You borrow \$1,000 from a bank to buy a car to use in your pizza delivery business.
	
	\end{enumerate}
	
\item Suppose GDP is \$8 trillion, taxes are \$1.5 trillion, private saving is \$0.5 trillion, and public saving is \$0.2 trillion. Assuming this economy is closed, calculate consumption, government purchases, national saving, and investment.	

\end{enumerate}

\section*{27.3. The Market for Loanable Funds}

\begin{enumerate}

\item The government institutes an investment tax credit -- a tax reduction for any firm that buys capital goods. What effect does the policy have on the market for loanable funds? Please use a supply and demand diagram in your response.

\item The government starts with a balanced budget, but then increases its spending so that it's running a budget deficit.

	\begin{enumerate}
	
	\item What effect does the budget deficit have on national saving?
	
	\item What effect does the policy have on the market for loanable funds? Please use a supply and demand diagram in your response. 
	
	\item You should have found that an increased budget deficit decreased investment. We call this phenomenon \underline{crowding out} because when the government borrows to finance its deficit, it crowds out private borrowers financing investment. The Republican party used to cite crowding out regularly when they advocated for spending cuts.
	
	\end{enumerate}
	
\item If a popular TV show on personal finance convinces Americans to save more for retirement, the \underline{\phantom{xxxxxxxxxx}} curve for loanable funds would shift, driving the equilibrium interest rate \underline{\phantom{xxxxxxxxxx}}.

	\begin{enumerate}
	
	\item supply; up
	
	\item supply; down
	
	\item demand; up
	
	\item demand; down
	
	\end{enumerate}
	
\item If the business community becomes more optimistic about the profitability of capital, the \underline{\phantom{xxxxxxxxxx}} curve for loanable funds would shift, driving the equilibrium interest rate \underline{\phantom{xxxxxxxxxx}}.

	\begin{enumerate}
	
	\item supply; up
	
	\item supply; down
	
	\item demand; up
	
	\item demand; down
	
	\end{enumerate}
	
\item Which of the following policy actions would unambiguously reduce the supply of loanable funds and crowd out investment?

	\begin{enumerate}
	
	\item an increase in taxes and a decrease in government spending
	
	\item a decrease in taxes together with an increase in government spending
	
	\item an increase in both taxes and government spending
	
	\item a decrease in both taxes and government spending
	
	\end{enumerate}

\end{enumerate}

\end{document}